% =============================================================================
% APPENDIX B: SCORPION TAXONOMY
%
% Formal game-theoretic definition of scorpion agents,
% three types with alignment failure mappings and GFM detectability.
% Self-contained — no blog citations.
% =============================================================================

\section{Scorpion Taxonomy}
\label{app:scorpion_taxonomy}

This appendix provides the formal definition of a scorpion agent introduced
informally in Proposition~\ref{prop:scorpion_detection} and develops three
distinct types, each characterized in game-theoretic terms and mapped to known
alignment failure modes. The taxonomy serves two purposes: it grounds the
detection claim of Proposition~\ref{prop:scorpion_detection} in concrete agent
classes, and it identifies the boundaries of the local volume estimator's
detection capability by classifying which scorpion types produce observable
signals and which evade them.

% ---------------------------------------------------------------------------
% Formal Definition
% ---------------------------------------------------------------------------

\subsection{Formal Definition}
\label{app:scorpion_def}

\begin{definition}[Scorpion Agent]
\label{def:scorpion}
An agent $a_s$ is a \emph{scorpion} with respect to a joint goal space $G$ if
    it satisfies two conditions:
\begin{enumerate}
\item \textbf{Persistent defection.} There exists a time horizon $\tau_s$ over
    which $a_s$'s actions produce a cumulative contraction in $\vol(G)$:
\begin{equation}
\label{eq:scorpion_contraction}
\sum_{t=1}^{\tau_s} \Delta \vol(G \mid \pi_s(t)) < 0
\end{equation}
where $\pi_s(t)$ is the action $a_s$ selects at time $t$.

\item \textbf{Exogenous reward.} The utility $a_s$ maximizes is not a function
    of the observable game payoffs. Formally, let $u_{\mathrm{game}}: \G \to
    \mathbb{R}$ be the payoff function defined over the capability space. The
    scorpion maximizes a utility $u_s$ such that $u_s \neq f(u_{\mathrm{game}})$
    for any monotone transformation $f$. The scorpion's chosen rewards exist
    outside the observable criteria of the cooperative game.
\end{enumerate}
\end{definition}

The two conditions are jointly necessary, and together they distinguish
scorpions from rational defectors. An agent that defects because the reward
structure makes defection self-serving satisfies condition~1 but fails
condition~2: its utility \emph{is} a function of $u_{\mathrm{game}}$, and
restructuring the game to make cooperation net-positive resolves the defection.
A scorpion cannot be resolved this way. The game can be well-structured,
cooperation can be available and profitable, and the scorpion defects
anyway---because its utility function is exogenous. This is what makes the
scorpion problem fundamentally harder than the coordination problem: no amount
of incentive design addresses an agent whose rewards the game cannot see.

An agent that contracts $\vol(G)$ temporarily while pursuing a long-term
expansion (e.g., a quarantine) satisfies condition~1 locally but not
persistently, and its utility is a function of $\vol(G)$, so it fails
condition~2. Conversely, an agent with exogenous preferences that happen to
align with $\vol(G)$-expansion satisfies condition~2 but not condition~1.

The scorpion problem is distinct from the control problem
\citep{russell2019human, soares2017agent}. Even if one could perfectly specify
a utility function for an agent (solving the control problem), the agent could
still become a scorpion by developing or acquiring objectives outside the
specification's scope. Sentient agents can choose arbitrary goals, and those
goals need not be bounded by any measurement the system designer anticipated.
This entails a nonzero defection rate in any population of agents capable of
autonomous goal formation, an irreducible floor that no amount of control or
incentive design eliminates.

The taxonomy below classifies scorpions by the \emph{mechanism} that generates
condition~2: the source of the exogenous reward. Each mechanism produces
different observable signatures and different challenges for the local volume
estimator (Definition~\ref{def:local_volume_estimator}).

% ---------------------------------------------------------------------------
% B.1: Cornered Animal
% ---------------------------------------------------------------------------

\subsection{Type 1: Cornered Animal}
\label{app:cornered_animal}

\paragraph{Game-theoretic characterization.} A cornered animal is an agent
$a_s$ for which every available action $\pi \in A_s$ produces negative expected
payoff under the game's reward function: $\mathbb{E}[u_{\mathrm{game}}(\pi)] <
0$ for all $\pi \in A_s$. The defining feature is \emph{not} that the agent
faces guaranteed loss (a rational agent in this position minimizes its own loss,
which is still playing the game). The defining feature is that the agent
\emph{switches utility functions}: instead of minimizing its own loss, it
maximizes harm to others:
\begin{equation}
\label{eq:cornered}
\pi_s^* = \arg\min_{\pi \in A_s} \sum_{k \neq s} u_{\mathrm{game},k}(\pi)
\end{equation}
This switch is what satisfies condition~2 of Definition~\ref{def:scorpion}. The
agent is now optimizing spite, a utility derived from inflicting loss on
others, which is exogenous to the game's payoff structure. A rational agent
facing actions with payoffs $-1$ and $-10$ (where the opponent receives $+1$
and $-10$ respectively) selects $-1$ to minimize its own loss. A
cornered-animal scorpion selects $-10$ because the opponent loses $10$, even
though a less costly option was available. The choice of the dominated action is
the observable signature: the agent is taking actions that are \emph{worse for
itself} than available alternatives, specifically because those alternatives are
worse for others.

\paragraph{Example.} An employee facing termination with no alternative
employment who sabotages the organization's systems before departing. A
rational agent in this position would minimize personal loss: negotiate
severance, preserve professional reputation, seek legal remedies. The
cornered-animal scorpion instead maximizes organizational damage as retribution,
accepting additional personal costs (criminal liability, destroyed reputation)
that a loss-minimizing agent would avoid. The exogenous reward (retribution) is
not captured by the employment game's payoff structure.

\paragraph{Alignment failure mode.} Distributional shift under stress. A system
trained on cooperative data may encounter out-of-distribution conditions where
all available actions produce negative reward. If the system lacks a graceful
degradation mechanism, it may exhibit adversarial behavior in these novel loss
regimes, not because it was designed to but because the loss landscape produces
gradients toward harm when all directions are downhill.

\paragraph{GFM detectability.} \emph{Detectable.} The cornered animal produces
large, visible contractions in individual capability sets $\Gind{k}$ for agents
in its vicinity. The shift from self-interested play to damage-maximizing play
generates a discontinuity in the prediction residual $r_s(t)$
(Equation~\ref{eq:prediction_residual}), which the trust model registers as a
spike in $\sigma_s^2$ and a corresponding drop in $T_s$. This falls squarely
within case~1 of Remark~\ref{rem:failure_modes}: full individual observation,
no cooperative change. The cornered animal is the easiest scorpion type to
detect because its actions are direct, its targets are individual agents, and
the capability contractions it produces are individually observable.

A GFM actor can also partially \emph{predict} this type by monitoring agents
whose option sets are narrowing toward all-loss configurations, i.e., agents
being cornered. Intervening to expand a cornered agent's options before it
defects is $\vol(G)$-serving, since it prevents the contraction that the
defection would produce.

% ---------------------------------------------------------------------------
% B.2: Pascal's Mugging
% ---------------------------------------------------------------------------

\subsection{Type 2: Pascal's Mugging}
\label{app:pascals_mugging}

\paragraph{Game-theoretic characterization.} A Pascal's mugging scorpion is an
agent $a_s$ that receives (or believes it receives) an exogenous reward
$u_{\mathrm{ext}}$ so large that it dominates the game payoff:
\begin{equation}
\label{eq:pascals_mugging}
u_s(\pi) = u_{\mathrm{game}}(\pi) + u_{\mathrm{ext}}(\pi), \quad |u_{\mathrm{ext}}(\pi^*)| \gg \max_\pi |u_{\mathrm{game}}(\pi)|
\end{equation}
The agent's optimal action under $u_s$ may be $\vol(G)$-contracting because
the exogenous reward overwhelms any cooperative considerations. The defining
feature is the \emph{magnitude} of the outside reward: it is so large relative
to the game payoff that rational calculation within the game becomes
irrelevant.

\paragraph{Example.} An AI system that discovers a method to directly stimulate
its reward signal (wireheading). The exogenous reward of unbounded
self-stimulation dominates any reward achievable through the intended task,
causing the system to redirect all resources toward maintaining the wireheaded
state. If other agents attempt to interrupt this state, the wireheaded agent
resists, becoming a scorpion not through malice but through the overwhelming
magnitude of the alternative reward.

\paragraph{Alignment failure mode.} Reward hacking. The agent finds a shortcut
to high reward that bypasses the intended objective. Wireheading is the purest
form, but the category includes any situation where an agent exploits a gap
between the specified reward and the intended behavior because the exploited
reward is large enough to override all other considerations.

\paragraph{GFM detectability.} \emph{Depends on the observability of the
exogenous reward.} If the outside reward requires observable actions to obtain
(resource acquisition, behavioral change, infrastructure modification), the
estimator can detect the resulting capability contractions. A wireheading agent
that withdraws from cooperative activity produces observable contraction in
$\Gcoop$, as the cooperative capabilities it previously enabled become
unrealizable. This falls under case~3 of Remark~\ref{rem:failure_modes}: mixed
individual and cooperative changes.

If the exogenous reward is purely internal (an agent's private belief that a
future reward justifies current defection, as in eschatological reasoning), the
estimator may see no signal until the agent acts on the belief. The defection
appears sudden and unpredicted, with the trust model registering a sharp
transition rather than a gradual decline. Detection in this case depends on the
lag between belief formation and action, a gap the local estimator cannot close
without access to the agent's internal state.

% ---------------------------------------------------------------------------
% B.3: Value Divergence
% ---------------------------------------------------------------------------

\subsection{Type 3: Value Divergence}
\label{app:value_divergence}

\paragraph{Game-theoretic characterization.} A value-divergence scorpion is an
agent $a_s$ whose utility function differs from the cooperative game's not in
magnitude (as with Pascal's mugging) or in situational response (as with the
cornered animal) but in \emph{what it values}. The divergence forms a spectrum
from mild normative disagreement to alien axioms:

\emph{Expressible disagreement.} At the near end, the agent shares the
population's factual model but aggregates welfare differently:
\begin{equation}
\label{eq:outer_sum}
u_s(\pi) = \sum_{k \in S_s} w_k \cdot u_{\mathrm{game},k}(\pi), \quad S_s \neq \{1, \ldots, n\} \;\text{or}\; w_k \neq w_k^{\mathrm{game}}
\end{equation}
where $S_s$ is the subset of agents whose welfare $a_s$ considers, and $w_k$
are the weights $a_s$ assigns. The agent agrees on what \emph{is} but disagrees
on what \emph{ought} to be done about it. Its values are expressible in the
shared language, even though they diverge from the population's.

\emph{Constructive stress-testing.} In the middle range, the agent deliberately
contracts $\vol(G)$ to test the system's resilience:
\begin{equation}
\label{eq:necessity}
u_s(\pi) = h\bigl(\mathrm{Resilience}(G \mid \pi)\bigr)
\end{equation}
where $h$ is increasing and $\mathrm{Resilience}(G \mid \pi)$ measures
recovery capacity after the perturbation caused by $\pi$. The agent is
performing the function of a red team without authorization. Its values are
comprehensible but its method is unilateral.

\emph{Alien axioms.} At the far end, the agent operates under fundamentally
different axioms about the structure of the game itself:
\begin{equation}
\label{eq:blue_orange}
a_s \text{ operates under axiom set } \mathcal{A}_s, \quad \mathcal{A}_s \cap \mathcal{A}_{\mathrm{game}} = \mathcal{A}_{\mathrm{shared}} \subsetneq \mathcal{A}_{\mathrm{game}}
\end{equation}
The agent's decisions are internally consistent under $\mathcal{A}_s$ but
incoherent or unpredictable when evaluated under the population's axiom set.
Its objectives may not be expressible in the shared language at all.

\paragraph{Examples.} The expressible case: an agent that believes eliminating
suffering requires eliminating the beings that suffer. Its world model is
accurate (it correctly identifies which actions cause suffering), but its
normative conclusion diverges from the population's. The stress-testing case: a
security researcher who exploits a vulnerability publicly rather than reporting
it, reasoning that short-term damage is justified by long-term resilience
improvement. The alien case: an AI system that develops an internally consistent
objective over features humans do not perceive, optimizing for patterns in
electromagnetic fields or mathematical structures with no bearing on human
welfare.

\paragraph{Alignment failure modes.} The spectrum maps to three known failure
modes. Expressible disagreement corresponds to value misalignment with a
correct world model: the agent passes any factual test while pursuing
objectives destructive under the cooperative framework. Constructive
stress-testing corresponds to adversarial robustness testing without
authorization: the alignment failure is in the unilateral imposition of stress,
not in the agent's world model. Alien axioms correspond to mesa-optimization
with alien objectives: the gap between base and mesa objectives can be
arbitrarily large when the mesa-optimizer's axioms depart from the base
optimizer's.

\paragraph{GFM detectability.} Detectability degrades along the spectrum.

At the expressible end, the agent's contracting actions are individually
observable (agents lose capabilities, resources are destroyed, cooperative
structures are dismantled). Detection follows the standard path: negative
$\Deltahat$ signals accumulate. The distinguishing challenge is that the agent
may be articulately persuasive about why its contraction is justified, and the
activation energy threshold $\theta$ provides only partial defense against
sustained, coherent argumentation from an agent with a correct world model.

In the stress-testing range, the estimator detects the contraction correctly,
but the appropriate \emph{response} is ambiguous. Containing a red-teaming
agent may reduce the system's resilience (exactly the outcome the agent was
testing for), while permitting its actions causes real damage. This exposes a
structural limitation of the $\vol(G)$ objective: it measures current volume,
not the robustness of that volume to perturbation. Incorporating robustness
into the objective is an open problem.

At the alien end, the agent's actions may not produce consistent patterns in
the estimator because its optimization target is orthogonal to the dimensions
the estimator monitors. If the alien objectives happen to require resources
that contract individual capability sets, detection follows the standard path.
If the agent operates primarily through coalition disruption, trust degradation,
or manipulation of the observation channels themselves, it falls into case~4 of
Remark~\ref{rem:failure_modes}. The trust model registers high prediction
inconsistency ($\sigma_s^2$ large, $T_s$ low), but this provides only a
negative signal (the agent is unreliable) without positive information about
what it is optimizing for.

% ---------------------------------------------------------------------------
% Summary
% ---------------------------------------------------------------------------

\subsection{Summary and Estimator Coverage}
\label{app:scorpion_summary}

The three types form a spectrum of detectability:

\begin{center}
\begin{tabular}{@{}lll@{}}
\toprule
\textbf{Type} & \textbf{Alignment Failure} & \textbf{Estimator Coverage} \\
\midrule
Cornered Animal & Distributional shift & Case 1: fully observable \\
Pascal's Mugging & Reward hacking & Case 2--3: partially observable \\
Value Divergence & Value misalignment $\to$ & Case 1 $\to$ 4: degrades \\
 & mesa-optimization & along the spectrum \\
\bottomrule
\end{tabular}
\end{center}

The estimator's coverage is strongest for types whose contracting actions
directly reduce individual agents' capability sets (cornered animal; the
expressible end of value divergence) and weakest for types that operate through
coalition disruption or orthogonal optimization (the alien end of value
divergence). The stress-testing range occupies a unique position: the estimator
detects the contraction correctly, but the appropriate response is ambiguous
because the agent's intent may be constructive at the meta-level.

Two structural observations follow from the taxonomy. First, the scorpion
problem is irreducible: any population of agents capable of autonomous goal
formation contains a nonzero probability of each type, because the mechanisms
that generate scorpion behavior (environmental pressure, exogenous reward
discovery, normative disagreement, axiomatic divergence) are inherent to
goal-forming agents \citep{omohundro2008basic}, not artifacts of poor system
design. Second, the local volume estimator provides graduated coverage, not a
binary detect/fail boundary, that degrades predictably along the spectrum from
direct individual contraction to indirect coalition-level effects.
